\documentclass[UTF8,fontset=none,aspectratio=169,10pt]{beamer}
\usepackage{ads-beamer-theme}
\title{高级分布式系统}
\subtitle{第六讲：分布式算法实验、拓扑与复杂度}
\author{陈鹏飞\quad 计算机学院}
\institute{中山大学}
\date{2025--2026 学年}

\begin{document}
\begin{frame}[plain]
\begin{tikzpicture}[remember picture,overlay]
  \fill[ADSBlue] (current page.north west) rectangle ([xshift=2.4cm]current page.south west);
  \fill[ADSOrange] ([xshift=2.4cm]current page.north west) rectangle ([xshift=2.55cm]current page.south west);
  \node[anchor=west,align=left] at ([xshift=3.3cm,yshift=1.0cm]current page.west) {\fontsize{32}{38}\selectfont\color{ADSBlue}\bfseries 高级分布式系统};
  \node[anchor=west,align=left] at ([xshift=3.3cm,yshift=-0.4cm]current page.west) {\fontsize{21}{27}\selectfont\color{ADSDarkBlue}第六讲\quad 算法实验与复杂度};
  \node[anchor=west,align=left] at ([xshift=3.3cm,yshift=-2.2cm]current page.west) {\large\color{ADSGray}把“会背算法”变成“会建模、会测量、会解释”};
  \node[anchor=east] at ([xshift=-0.8cm,yshift=-3.2cm]current page.east) {\small\color{ADSGray}陈鹏飞\quad 中山大学\quad 2025--2026};
  \draw[ADSBlue,line width=0.8pt] ([xshift=3.3cm,yshift=-1.35cm]current page.west) -- ([xshift=12.6cm,yshift=-1.35cm]current page.west);
\end{tikzpicture}
\end{frame}

\sectionframe{一、算法不是代码片段}{算法的输入、系统模型和停止条件决定了“正确”到底是什么意思。}

\begin{frame}{本讲学习目标}
\begin{enumerate}
  \item 区分物理拓扑、逻辑拓扑与叠加拓扑；
  \item 理解控制算法与应用执行如何并行存在；
  \item 复习集中式/分布式、对称/非对称、确定性/非确定性算法；
  \item 掌握同步和异步生成树实验的状态、消息和复杂度；
  \item 能把算法复杂度拆成节点空间、消息数、消息大小和时间。
\end{enumerate}
\vfill
\callout{实验主线}{先画事件和状态，再写伪代码；先定义测量口径，再比较优化。}
\end{frame}

\begin{frame}{物理拓扑、逻辑拓扑、覆盖拓扑}
\begin{center}
\begin{tikzpicture}[>=Latex,node distance=12mm and 18mm]
\nodebox[minimum width=2.5cm]{ADSGray}{phy}{物理网络\\路由器/链路};
\nodebox[right=of phy,minimum width=2.5cm]{ADSBlue}{logic}{逻辑信道\\进程视图};
\nodebox[right=of logic,minimum width=2.5cm]{ADSOrange}{over}{覆盖图\\环/树/网格};
\draw[\arrow,ADSGray] (phy)--node[above,font=\scriptsize]{抽象}(logic);
\draw[\arrow,ADSGray] (logic)--node[above,font=\scriptsize]{叠加}(over);
\end{tikzpicture}
\end{center}
\begin{itemize}
  \item 物理邻居不一定是算法中的逻辑邻居；
  \item 覆盖图的目标是高效广播、搜索、路由或收集信息；
  \item 实验必须说明测的是哪个层次的边和延迟。
\end{itemize}
\end{frame}

\begin{frame}{应用执行与控制算法}
\begin{columns}[T,onlytextwidth]
\column{.48\linewidth}
\blocktitle{应用执行}
\begin{itemize}
  \item 处理业务消息和数据；
  \item 不能因为控制协议而丢失语义；
  \item 例如训练、查询、事务。
\end{itemize}
\column{.48\linewidth}
\blocktitle{控制算法}
\begin{itemize}
  \item 监控、建树、检测全局状态；
  \item 叠加在应用上，事件要尽量透明；
  \item 例如快照、共识、故障检测。
\end{itemize}
\end{columns}
\vfill
\callout{解释}{控制协议的消息也消耗带宽和队列空间，因此“透明”不等于“没有成本”。}
\end{frame}

\begin{frame}{集中式与分布式算法}
\begin{tabularx}{\linewidth}{>{\bfseries}p{3cm}XX}
\toprule
集中式 & 角色不对称，控制集中在一个服务或领导者。 & 易实现，但有瓶颈和单点故障。 \\
分布式 & 多个节点共同计算，状态和通信分散。 & 更有韧性，但状态交错与证明更复杂。 \\
对称/非对称 & 节点是否执行同一规则。 & 对称更易扩展，非对称常用特殊节点打破对称。 \\
\bottomrule
\end{tabularx}
\end{frame}

\begin{frame}{匿名、uniform 与 adaptive}
\begin{itemize}
  \item 匿名算法不使用进程 ID 作运行时决策，结构优雅但可能无法选 leader；
  \item uniform 算法不把系统规模 n 写死，适合节点动态加入和离开；
  \item adaptive 算法按当前参与者数 k 分析复杂度，而不是最大规模 n；
  \item 选择模型时要问：算法是否知道 n、diameter、根节点或故障上界？
\end{itemize}
\end{frame}

\begin{frame}{确定性与非确定性接收}
\begin{itemize}
  \item 确定性 receive 指定消息来源，重放实验更容易复现；
  \item 非确定性 receive 可以从任意来源取消息，执行可能产生不同事件偏序；
  \item 异步网络中的排队、CPU 调度和拥塞会改变交错；
  \item 调试时记录消息来源、逻辑时间和随机种子，才能重演问题。
\end{itemize}
\end{frame}

\sectionframe{二、系统模型与故障模型}{没有模型边界，复杂度和正确性都无法解释。}

\begin{frame}{同步与异步系统}
\begin{tabularx}{\linewidth}{>{\bfseries}p{3cm}XX}
\toprule
同步 & 消息延迟、时钟漂移、一步执行时间都有已知上界。 & 可用轮次和超时推理。 \\
异步 & 上述上界不存在或未知。 & 只能依靠因果、确认和故障检测猜测进展。 \\
部分同步 & 最终存在上界，但系统不知道何时开始。 & 工程上常见，协议需容忍早期误判。 \\
\bottomrule
\end{tabularx}
\end{frame}

\begin{frame}{进程故障从轻到重}
\begin{enumerate}
  \item fail-stop：停止后其他节点能得知；
  \item crash：停止但其他节点只能猜测；
  \item receive omission：收不到部分消息；
  \item send omission：发不出应发送的部分消息；
  \item general omission：发送和接收都可能遗漏；
  \item Byzantine：任意行为，甚至伪造和篡改消息。
\end{enumerate}
\end{frame}

\begin{frame}{链路故障与认证}
\begin{itemize}
  \item 链路 crash：链路完全停止；
  \item 链路 omission：部分消息丢失；
  \item 链路 Byzantine：制造假消息或修改内容；
  \item 认证可以让“某故障进程伪造另一个正确进程的消息”变得可检测；
  \item 但认证不能阻止拒绝服务，仍需冗余和限流。
\end{itemize}
\end{frame}

\begin{frame}{执行抑制（freezing）}
\begin{itemize}
  \item inhibitory protocol 会暂停某些正常事件，等待控制操作完成；
  \item non-inhibitory：任何执行中都不禁用事件；
  \item locally inhibitory：延迟由本地状态决定；
  \item globally inhibitory：某事件是否被延迟取决于远端 receive；
  \item 实验报告必须说明控制协议是否改变应用吞吐和尾延迟。
\end{itemize}
\end{frame}

\begin{frame}{Wait-free 与通信信道}
\begin{itemize}
  \item wait-free 操作在最多 $n-1$ 个进程失败时仍须在有限步完成；
  \item 鲁棒但通常更昂贵，某些生产者—消费者问题可能无法 wait-free；
  \item 点到点信道可以 FIFO 或非 FIFO；
  \item 应用层“看起来 FIFO”不代表跨多跳路径或重试仍保持 FIFO。
\end{itemize}
\end{frame}

\sectionframe{三、复杂度：不要只报一个 O()}{消息数量、消息大小、并行度和恢复成本共同决定系统代价。}

\begin{frame}{复杂度指标}
\begin{itemize}
  \item 节点本地空间与全局空间；
  \item 节点本地时间与系统总时间；
  \item 消息数、每条消息大小、消息 hop 和并发度；
  \item send/receive 事件、广播次数和同步操作数；
  \item 共享内存算法还要统计内存大小、读写和同步次数。
\end{itemize}
\vfill
\callout{实验提醒}{两种算法都写 O(m)，但 m 是边数、消息数还是 payload 字节？不先定义，数字没有可比性。}
\end{frame}

\begin{frame}{消息复杂度的三个层次}
\begin{tabularx}{\linewidth}{>{\bfseries}p{3.2cm}XX}
\toprule
消息数 & 网络拥塞和控制开销。 & 一条大消息与一百条小消息不同。 \\
消息大小 & 传输时间、缓存和序列化成本。 & checkpoint/模型权重尤其明显。 \\
消息时间 & hops、并行发送、队列等待共同决定。 & 不能用 hop 数直接替代端到端延迟。 \\
\bottomrule
\end{tabularx}
\end{frame}

\begin{frame}{实验设计：变量与对照}
\begin{enumerate}
  \item 固定拓扑、消息大小和故障模型；
  \item 只改变一个因素：节点数、直径、负载或丢包率；
  \item 同时记录平均值、P95/P99 和恢复时间；
  \item 用基线算法和优化算法做对照；
  \item 把正确性检查和性能测量分开运行。
\end{enumerate}
\end{frame}

\sectionframe{四、同步单发起者生成树}{Flooding 的关键不是“发出去”，而是确定 parent、深度和终止。}

\begin{frame}{同步 1-initiator spanning tree}
\begin{itemize}
  \item 根节点 visited=1、depth=0，向所有邻居发送 QUERY；
  \item 每一轮未访问节点若收到 QUERY，选择一个发送者为 parent；
  \item 设置 depth=round，向除发送者外邻居继续扩散；
  \item 运行到网络直径后停止；
  \item 生成的树是按轮次发现的 BFS 风格树。
\end{itemize}
\end{frame}

\begin{frame}{同步生成树示例}
\begin{center}
\begin{tikzpicture}[>=Latex,node distance=13mm]
\nodebox[minimum width=1.7cm]{ADSBlue}{a}{A\,(0)};
\nodebox[right=of a,minimum width=1.7cm]{ADSOrange}{b}{B\,(1)};
\nodebox[right=of b,minimum width=1.7cm]{ADSGreen}{c}{C\,(2)};
\nodebox[below=of a,minimum width=1.7cm]{ADSGray}{f}{F\,(1)};
\nodebox[below=of b,minimum width=1.7cm]{ADSGray}{e}{E\,(2)};
\nodebox[below=of c,minimum width=1.7cm]{ADSGray}{d}{D\,(3)};
\draw[ADSBlue,very thick] (a)--(b)--(c);
\draw[ADSBlue,very thick] (a)--(f)--(e)--(d);
\draw[ADSGray,dashed] (b)--(e) (c)--(d);
\end{tikzpicture}
\end{center}
\begin{itemize}
  \item 圆括号表示发现轮次；蓝线是 parent 关系；
  \item 同一轮收到多个 QUERY 时任选一个会影响树形，但不影响可达性。
\end{itemize}
\end{frame}

\begin{frame}{同步算法的终止与复杂度}
\begin{itemize}
  \item 根知道直径上界时可等待 d 轮；
  \item 若节点只设置 parent 就停止，其他节点如何知道它已处理完？
  \item 每条边最多发送 1--2 条 QUERY，消息复杂度约 $[l,2l]$；
  \item 本地空间 $O(\text{degree})$，时间约为 degree + diameter。
\end{itemize}
\end{frame}

\sectionframe{五、异步生成树：确认、拒绝与局部终止}{异步系统没有轮次，必须用消息把“我是不是你的 child”说清楚。}

\begin{frame}{异步 1-initiator 的状态}
\begin{itemize}
  \item parent：第一次接受 QUERY 的邻居；
  \item Children：收到 ACCEPT 的邻居集合；
  \item Unrelated：收到 REJECT 的邻居集合；
  \item 只有当所有非 parent 邻居都被 ACCEPT 或 REJECT 覆盖，节点才能终止。
\end{itemize}
\end{frame}

\begin{frame}{异步消息处理}
\begin{enumerate}
  \item 根向所有邻居发送 QUERY；
  \item 节点第一次收到 QUERY，设置 parent，回 ACCEPT，并向其他邻居继续 QUERY；
  \item 后续 QUERY 来自非 parent，回 REJECT；
  \item 收到 ACCEPT 加入 Children，收到 REJECT 加入 Unrelated；
  \item 本地覆盖完邻居后终止或向 parent 回报。
\end{enumerate}
\end{frame}

\begin{frame}{为什么需要 REJECT？}
\begin{itemize}
  \item 没有 REJECT，发送 QUERY 的节点不知道邻居是否已被别人吸收；
  \item 终止条件无法区分“没有 child”和“还没回复”；
  \item 代价是每条非树边需要额外的否定确认；
  \item 可以用显式邻居状态、序列号或聚合 ACK 替代，但复杂度不会凭空消失。
\end{itemize}
\end{frame}

\begin{frame}{异步生成树的复杂度}
\begin{itemize}
  \item 每条树边约 2 条消息（QUERY + ACCEPT）；
  \item 每条非树边约 2 条消息（QUERY + REJECT）；
  \item 总消息约 $[2l,4l]$；
  \item 没有同步轮次保证，树高最坏可达 $n-1$；
  \item 时间复杂度按消息 hop 和最长因果链测量。
\end{itemize}
\end{frame}

\sectionframe{六、并发发起者与拓扑算法}{没有预先指定 root 时，算法必须合并、抑制或比较多个实例。}

\begin{frame}{并发发起者的三种选择}
\begin{enumerate}
  \item 合并多个部分生成树：需要额外协议，容易形成环；
  \item 允许一个实例继续，其他实例被抑制；
  \item 按 root ID 选择优先级，更大的 root 取代更小的 root。
\end{enumerate}
\vfill
\callout{关键点}{“只保留一个实例”需要传播 root 身份、清理旧状态，并证明旧实例不会重新复活。}
\end{frame}

\begin{frame}{并发生成树的状态}
\begin{itemize}
  \item myroot：当前节点认为的根；
  \item QUERY(root)：带根身份传播；
  \item 新 root 更大时清空旧 Children/Unrelated 并切换 parent；
  \item 新 root 相同才处理 ACCEPT/REJECT；旧 root 消息被忽略。
\end{itemize}
\end{frame}

\begin{frame}{树上的 broadcast 与 convergecast}
\begin{center}
\begin{tikzpicture}[>=Latex,node distance=14mm]
\nodebox[minimum width=1.8cm]{ADSBlue}{rootx}{root};
\nodebox[below left=of rootx,minimum width=1.8cm]{ADSOrange}{l1}{child};
\nodebox[below right=of rootx,minimum width=1.8cm]{ADSOrange}{r1}{child};
\nodebox[below=of l1,minimum width=1.8cm]{ADSGreen}{l2}{leaf};
\nodebox[below=of r1,minimum width=1.8cm]{ADSGreen}{r2}{leaf};
\draw[\arrow,ADSBlue] (rootx)--(l1) (rootx)--(r1);
\draw[\arrow,ADSBlue] (l1)--(l2) (r1)--(r2);
\draw[\arrow,ADSOrange] (l2)--(l1) (r2)--(r1) (l1)--(rootx) (r1)--(rootx);
\end{tikzpicture}
\end{center}
\begin{itemize}
  \item broadcast 从 root 向下，convergecast 从叶子向上；
  \item 同一棵树可复用来收集统计、汇总 checkpoint 或传播配置。
\end{itemize}
\end{frame}

\sectionframe{七、AI 基础设施实验}{用相同的模型和指标，验证控制算法是否真的改善系统。}

\begin{frame}{实验一：拓扑对模型分发延迟的影响}
\begin{enumerate}
  \item 在 ring、tree、mesh 三种逻辑拓扑中放置相同数量节点；
  \item 让一个 2 GB 模型从根节点广播；
  \item 记录完成时间、峰值带宽、重复消息和失败恢复时间；
  \item 解释：树的 hop 少但根可能成为瓶颈，mesh 复制多但更抗故障。
\end{enumerate}
\end{frame}

\begin{frame}{实验二：故障模型对 checkpoint 收集的影响}
\begin{itemize}
  \item fail-stop：收集器可以快速排除节点；
  \item crash：需要 heartbeat 和超时，存在误判；
  \item omission：重试和幂等 ID 防止重复提交；
  \item Byzantine：必须验证摘要、签名和来源；
  \item 输出不只报平均时间，还要给出丢失比例和尾延迟。
\end{itemize}
\end{frame}

\begin{frame}{实验三：公平性与自适应复杂度}
\begin{itemize}
  \item 改变当前竞争 GPU 的作业数 k，但保持集群规模 n 不变；
  \item 测量每个作业获得资源的等待时间和消息数；
  \item 如果复杂度随 k 而非 n 增长，说明算法具有 adaptive 特征；
  \item 检查最慢作业是否长期饥饿。
\end{itemize}
\end{frame}

\sectionframe{八、总结与实验作业}{把课件中的伪代码，转化为可观察、可复现、可质疑的系统。}

\begin{frame}{算法实验报告模板}
\begin{enumerate}
  \item 系统模型：节点、信道、时钟、故障；
  \item 算法状态：每个变量的含义、不变量和更新条件；
  \item 消息协议：发送方、接收方、字段、重试与幂等；
  \item 正确性：安全、活性、终止或收敛；
  \item 复杂度：空间、消息、时间、payload、恢复；
  \item 实验：参数、基线、曲线、失败案例和解释。
\end{enumerate}
\end{frame}

\begin{frame}{本讲知识地图}
\begin{center}
\begin{tikzpicture}[>=Latex,node distance=11mm and 16mm]
\nodebox[minimum width=2.3cm]{ADSBlue}{topo}{拓扑抽象};
\nodebox[right=of topo,minimum width=2.3cm]{ADSOrange}{modelx}{系统/故障模型};
\nodebox[below=of topo,minimum width=2.3cm]{ADSGreen}{tree}{生成树};
\nodebox[right=of tree,minimum width=2.3cm]{ADSGray}{metric}{复杂度指标};
\nodebox[below=of metric,minimum width=2.5cm]{ADSBlue}{ai6}{AI 实验};
\draw[\arrow,ADSGray] (topo)--(modelx);
\draw[\arrow,ADSGray] (topo)--(tree);
\draw[\arrow,ADSGray] (modelx)--(metric);
\draw[\arrow,ADSGray] (tree)--(metric);
\draw[\arrow,ADSGray] (metric)--(ai6);
\end{tikzpicture}
\end{center}
\callout{带走一句话}{算法的价值不只在于“能运行”，而在于知道它依赖什么、保证什么、花费什么。}
\end{frame}

\begin{frame}{课后作业}
\begin{enumerate}
  \item 实现同步和异步 1-init spanning tree，注入随机消息延迟；
  \item 对比 QUERY/ACCEPT/REJECT 的消息数量和终止时间；
  \item 增加一个崩溃节点，观察树是否能被检测和修复；
  \item 将实验结果映射到模型分发或 checkpoint 收集场景。
\end{enumerate}
\vfill
\smallnote{课后可复现实验：比较同步/异步生成树的消息、延迟和恢复成本。}
\end{frame}

\end{document}
